\itemhead{1.2 例题 · 例题9（教材p.22）}{XJ-B1-C01-EX-P022-Y0891}
\label{XJ-B1-C01-EX-P022-Y0891}
{\footnotesize 湘教版教材p.22 / PDF第29页 / 1.2 例题 / 例题9}
\field{原题}{写出下列含量词命题的否定：
\subq{1}$p$：所有二次函数的图象都是轴对称图形；
\subq{2}$q$：$\exists x\in\mathbb R,\ -2x^2+4x-3>0$；
\subq{3}$r$：三角形的垂心都在其外部；
\subq{4}$s$：有一个小于 210 的正整数至少有 4 个质因数。}
\field{参考答案}{（1）存在一个二次函数，其图象不是轴对称图形；（2）$\forall x\in\mathbb R,-2x^2+4x-3\le0$；（3）存在一个三角形，其垂心不在外部；（4）每一个小于210的正整数的质因数都少于4个。}
{\small \field{关联知识}{L04 含一个量词命题的否定（教材帮p.39、40 / PDF48、49）}
\field{方法与易错点}{分三步书写：保留论域、交换量词、否定断言，再判断新命题真假。 易错：只改量词不改结论；把不等号否定成反向但遗漏等号。}
\field{答案依据}{例题原解位置：湘教版教材p.22及题库原图；独立编写可编辑参考解答。答案册此部分主要收练习、习题，不伪造例题答案册位置。}}
